\bookchapter{Bab 4: Memasang dan Mengelola Aplikasi — “Paket Keterampilan” Ponsel}

\chapterintro{Aplikasi membuat ponsel bisa melakukan apa saja, tetapi juga perlu dikelola dengan baik. Sama seperti Anda memberikan meja kerja dan mengalokasikan izin kepada karyawan baru, memasang aplikasi juga memerlukan proses standar.}

\booksection{Toko Resmi: “Supermarket Aplikasi” Paling Aman}

\begin{itemize}
  \item \textbf{Android}: Google Play (luar negeri) atau toko bawaan masing-masing merek (seperti Huawei AppGallery, Mi Store).
  \item \textbf{: Pengguna biasa tidak dapat mengaktifkan mode pengembang lengkap, membutuhkan akun Apple Developer berbayar.}: Hanya App Store.
  \item \textbf{HarmonyOS}: Huawei AppGallery.
\end{itemize}

Arti penting keamanan inti: telah melalui peninjauan, pembaruan terjamin, pembayaran aman. Panduan tindakan:\textbf{prioritaskan hanya mengunduh dari toko resmi!}

\booksection{Mode Murni/Proteksi: Aktifkan “Tembok Anti-Virus” Ponsel}

Merek domestik mainstream (Huawei, Xiaomi, OPPO, vivo, Honor) semua memiliki\textbf{Mode Murni}atau\textbf{Mode Proteksi}bawaan, ini adalah sakelar keamanan tingkat sistem.\textbf{Sangat disarankan agar semua pengguna mengaktifkannya}.

\booksubsection{Apa sebenarnya itu?}

Fungsi inti Mode Murni (juga disebut “Mode Proteksi”) adalah memblokir pemasangan aplikasi pihak ketiga\textbf{: saat aktif, sistem akan}hanya mengizinkan pemasangan aplikasi dari toko aplikasi resmi\textbf{, paket instalasi yang diunduh dari browser, QR code, perangkat lunak chat, dan saluran lain akan langsung diblokir, mencegah perangkat lunak berbahaya dari sumbernya.}2. Pengguna yang Sesuai

\booksubsection{Pengguna lanjut usia}

\begin{itemize}
  \item \textbf{: Tidak bisa membedakan iklan "resmi" dan "palsu", Mode Murni adalah jimat paling aman.}Pengguna yang tidak sensitif terhadap izin aplikasi
  \item \textbf{: Tidak perlu memasang aplikasi niche, beta, game server luar negeri.}Semua pengguna yang ingin meningkatkan keamanan perangkat
  \item \textbf{: Bahkan penggemar, lebih baik aktif daripada tidak — jika benar-benar butuh sementara, matikan sebentar saja.}3. Jalur Aktivasi (contoh Huawei/Xiaomi)
\end{itemize}

\booksubsection{Huawei/HarmonyOS}

\begin{itemize}
  \item \textbf{: [Pengaturan] > [Sistem \& Pembaruan] >}[Mode Murni]\textbf{, pilih "Aktifkan".}Xiaomi
  \item \textbf{Xiaomi}[Lindungi Privasi]\textbf{atau [Pengaturan] > [Pengaturan Aplikasi] > [Pemasangan Aplikasi], matikan "Pasang aplikasi tidak dikenal".}OPPO/vivo
  \item \textbf{: Sebagian model dapat di-unlock, perlu pengajuan.}[Pasang Aplikasi Sumber Eksternal]\textbf{, matikan otorisasi satu per satu.}4. Mekanisme Inspeksi Keamanan Toko Aplikasi
\end{itemize}

\booksubsection{Toko aplikasi bukan “pajangan”, di belakangnya ada jalur produksi keamanan lengkap:}

Pemindaian virus

\begin{itemize}
  \item \textbf{: Setiap aplikasi sebelum dipublikasikan akan dipindai bersama oleh beberapa mesin antivirus, mengidentifikasi karakteristik trojan dan virus yang diketahui.}Audit izin
  \item \textbf{: Toko akan meninjau apakah izin yang diminta aplikasi "masuk akal" — aplikasi senter meminta izin kontak? Langsung ditolak atau ditarik.}Sertifikasi pengembang
  \item \textbf{: Publikasi memerlukan kualifikasi perusahaan, KTP pribadi atau kartu kredit. Jika ada masalah bisa ditemukan orangnya, bukan "aplikasi hitam".}Pemantauan berkelanjutan
  \item \textbf{: Aplikasi yang sudah dipublikasikan juga akan terus dipindai ulang, perilaku berbahaya yang baru ditemukan akan dipaksa ditarik.}3. Sumber Pihak Ketiga dan "Parasit Elektronik": Risiko Sangat Tinggi!
\end{itemize}

\booksection{Sumber Pihak Ketiga dan “Parasit Elektronik”: Risiko Sangat Tinggi!}

“parasit elektronik”\textbf{.}.

\textbf{Merujuk pada}\par

perangkat lunak “versi crack” yang ditanam trojan\textbf{— Anda pikir “menghemat uang”, sebenarnya “terjebak”. Bentuk umum:}Perangkat lunak berbayar “versi crack”, “versi gratis”, “versi tanpa iklan”;

\begin{itemize}
  \item Alat abu-abu seperti "suka otomatis", "grabs red packet", "cheat game";
  \item Aplikasi "versi internal" yang direkomendasikan "iklan kecil" di grup chat.
  \item Program-program ini tampak bisa digunakan, di belakang mungkin
\end{itemize}

mencuri kontak, membaca SMS, menyadap mikrofon, menambang di latar belakang\textbf{. Mode Murni adalah}garis pertahanan pertama\textbf{melawan mereka — langsung tidak membiarkan Anda memasangnya.}Saran garis pertahanan kedua:

\booksubsection{Pemindaian virus penuh berkala}

\begin{itemize}
  \item \textbf{: Gunakan pusat keamanan bawaan ponsel (seperti Huawei Phone Manager, Mi Security Center) memindai sekali sebulan.}Berhati-hati memberikan izin sensitif
  \item \textbf{: Kontak, SMS, lokasi, kamera, mikrofon — kelima "izin berisiko tinggi" ini harus diaudit ketat.}Perhatikan anomali konsumsi daya dan kuota
  \item \textbf{: Ponsel tiba-tiba panas, baterai cepat berkurang, kuota berjalan tengah malam — kemungkinan besar ada "parasit" tersembunyi di latar belakang.}Jika terpaksa harus memasang aplikasi pihak ketiga:
\end{itemize}

Pastikan sumber benar-benar tepercaya.

\begin{enumerate}
  \item Periksa persyaratan izin dengan teliti sebelum memasang.
  \item Pindai dengan aplikasi keamanan setelahnya.
  \item Saran terakhir: Di pengaturan sistem, matikan izin "Pasang aplikasi tidak dikenal".
\end{enumerate}

4. Pasang, Perbarui, Copot

\booksection{Pasang, Perbarui, Copot}

\begin{itemize}
  \item \textbf{: Di toko klik "Dapatkan" atau "Pasang", setelah terpasang pertama kali dibuka akan meminta izin.}Pembaruan
  \item \textbf{:}:\textbf{Utamanya untuk memperbaiki kerentanan keamanan. Disarankan mengaktifkan pembaruan otomatis.}Copot
  \item \textbf{: Tahan lama ikon seret ke area copot. Untuk membersihkan lebih tuntas, buka [Pengaturan] > [Manajemen Aplikasi] cari dan copot.}5. Manajemen Izin: Anda "Bos" Aplikasi
\end{itemize}

\booksection{Manajemen Izin: Anda “Bos” Aplikasi}

Opsi izin

\begin{itemize}
  \item \textbf{: "Selalu izinkan", "Izinkan hanya saat digunakan", "Larang", "Tanya setiap kali".}Prinsip pengelolaan
  \item \textbf{:}:
  \begin{enumerate}
    \item Audit setelahnya: Periksa secara berkala di pengaturan izin siapa yang "melampaui wewenang".
    \item Kontekstual: Seperti aplikasi taksi, berikan izin lokasi hanya saat digunakan.
    \item Izin yang Perlu Diawasi Ketat
  \end{enumerate}
\end{itemize}

\begin{tcbwarn}{Kontak/SMS: Mungkin digunakan untuk mengirim SMS spam.}
\begin{itemize}
  \item Lokasi real-time: Mungkin dilacak pergerakannya.
  \item Kamera/mikrofon: Mungkin diam-diam mengambil foto/rekaman.
  \item Akses file: Mungkin membaca file pribadi Anda.
  \item 6. Komponen Paket Instalasi Perangkat Lunak
\end{itemize}

\end{tcbwarn}

\booksection{Pernahkah Anda bertanya-tanya, paket instalasi yang diunduh dari toko aplikasi, apa sebenarnya isinya? Memahami komposisi paket instalasi dapat membantu Anda lebih memahami cara kerja aplikasi, juga dapat membantu saat mengalami masalah mengetahui dari mana harus memeriksa.}

1. Struktur File APK Android

\booksubsection{APK (Android Package Kit) pada dasarnya adalah arsip ZIP, setelah diekstrak dapat dilihat direktori dan file inti berikut:}

AndroidManifest.xml

\begin{itemize}
  \item \textbf{: "KTP" dan "buku petunjuk" aplikasi. Mencatat nama paket aplikasi, nomor versi, izin yang diperlukan, versi sistem yang didukung, Activity masuk, dan informasi inti lainnya. Sistem membaca file ini pertama kali saat memasang aplikasi.}classes.dex
  \item \textbf{: Kode inti aplikasi. Kode sumber Java/Kotlin setelah dikompilasi menjadi bytecode Dalvik, disimpan dalam file ini. Satu aplikasi mungkin memiliki beberapa file dex (seperti classes2.dex, classes3.dex).}resources.arsc
  \item \textbf{: Tabel indeks sumber daya yang dikompilasi. Mencatat ID semua sumber daya dan korespondensinya.}Direktori res/
  \item \textbf{: Folder sumber daya. Berisi gambar, file layout, string, definisi warna, dll. Gambar dengan resolusi berbeda ditempatkan di folder drawable berbeda (seperti drawable-hdpi, drawable-xhdpi).}Direktori assets/
  \item \textbf{: Folder sumber daya asli. Menyimpan sumber daya yang tidak dikompilasi, seperti file halaman web, font, audio, dll. Saat aplikasi berjalan dapat dibaca melalui AssetManager.}Direktori lib/
  \item \textbf{: Folder pustaka native. Menyimpan file .so yang telah dikompilasi (hasil kompilasi kode C/C++), diklasifikasikan berdasarkan arsitektur CPU (seperti arm64-v8a, armeabi-v7a). Game dan aplikasi yang membutuhkan komputasi performa tinggi biasanya memiliki direktori ini.}Direktori META-INF/
  \item \textbf{: Folder informasi tanda tangan. Berisi tiga file MANIFEST.MF, CERT.SF, CERT.RSA, digunakan untuk memverifikasi integritas dan tanda tangan APK.}2. Struktur File IPA iOS
\end{itemize}

\booksubsection{IPA (iOS App Store Package) juga arsip terkompresi, strukturnya sedikit berbeda:}

Direktori Payload/

\begin{itemize}
  \item \textbf{: Lokasi konten inti. Di dalamnya ada folder .app, berisi semua konten aplikasi.}Info.plist
  \item \textbf{: File konfigurasi aplikasi, mirip dengan AndroidManifest.xml Android. Mencatat ID aplikasi, versi, persyaratan izin, perangkat yang didukung, dll.}File yang dapat dieksekusi
  \item \textbf{: Program utama dengan nama yang sama dengan folder .app, adalah file biner format Mach-O yang dikompilasi.}Direktori Frameworks/
  \item \textbf{: Pustaka dinamis dan kerangka kerja yang menjadi dependensi aplikasi.}File sumber daya
  \item \textbf{: Gambar, audio, video, dan sumber daya lainnya, biasanya disimpan di dalam folder .app.}Direktori \_CodeSignature/
  \item \textbf{: Informasi tanda tangan kode, digunakan untuk memverifikasi integritas dan asal aplikasi.}3. Mekanisme Verifikasi Tanda Tangan
\end{itemize}

\booksubsection{Mengapa aplikasi yang diunduh dari toko resmi lebih aman? Verifikasi tanda tangan adalah kuncinya.}

Tanda tangan Android

\begin{itemize}
  \item \textbf{: Pengembang menandatangani APK dengan kunci privat, sistem memverifikasi dengan kunci publik saat pemasangan. Jika APK diubah (misalnya ditanam virus), tanda tangan akan tidak valid, sistem akan menolak memasang. Google Play juga mewajibkan App Signing, Google menyimpan kunci tanda tangan, lebih lanjut melindungi keamanan aplikasi.}Tanda tangan iOS
  \item \textbf{: Peninjauan Apple lebih ketat. Setiap IPA harus ditandatangani sertifikat yang dikeluarkan Apple, jika tidak, tidak dapat dipasang. Inilah sebabnya iOS hanya dapat memasang aplikasi dari App Store atau sertifikat perusahaan. Aplikasi yang ditandatangani sertifikat perusahaan memiliki masa berlaku terbatas, dan Apple dapat mencabutnya kapan saja.}Fungsi tanda tangan
  \item \textbf{:}:
  \begin{enumerate}
    \item Verifikasi integritas: Memastikan aplikasi tidak diubah selama transmisi.
    \item Mencegah penyamaran: Aplikasi berbahaya tidak dapat menyamar sebagai aplikasi lain.
    \item 4. Ukuran Paket Instalasi dan Optimasi
  \end{enumerate}
\end{itemize}

\booksubsection{Anda mungkin menemukan, ukuran aplikasi yang ditampilkan di toko aplikasi berbeda dengan ruang yang digunakan setelah dipasang. Ini karena:}

Ukuran unduhan

\begin{itemize}
  \item \textbf{: Ukuran APK/IPA setelah dikompres.}Ukuran pemasangan
  \item \textbf{: Ukuran setelah diekstrak, biasanya jauh lebih besar dari ukuran unduhan.}Penggunaan runtime
  \item \textbf{: Saat aplikasi berjalan juga akan menghasilkan cache, file data, menempati lebih banyak ruang.}Sarana optimasi aplikasi modern:
\end{itemize}

App Bundle (Android)

\begin{itemize}
  \item \textbf{: Teknologi distribusi dinamis Google Play. Saat mengunduh, pengguna hanya mendapatkan sumber daya yang sesuai dengan perangkatnya (seperti hanya mengunduh gambar resolusi layar yang sesuai), sangat mengurangi ukuran unduhan.}App Thinning (iOS)
  \item \textbf{: Teknologi serupa, Apple akan otomatis memotong sumber daya yang tidak diperlukan sesuai perangkat.}Unduhan sesuai permintaan
  \item \textbf{: Aplikasi besar seperti game dapat mengunduh paket inti terlebih dahulu, kemudian mengunduh paket sumber daya.}Tahukah Anda: Mengapa ada paket instalasi aplikasi yang sangat besar?
\end{itemize}

\begin{tcbinfo}{Aplikasi jenis game biasanya paling besar, karena berisi banyak tekstur HD, model 3D, file audio. Aplikasi jenis sosial juga cukup besar, karena harus mendukung berbagai modul fungsi. Aplikasi jenis alat biasanya paling kecil. Jika Anda menemukan aplikasi alat tertentu sangat besar, waspadai apakah dibundel dengan fungsi yang tidak perlu.}
Praktik: Audit Izin 3 Aplikasi

\end{tcbinfo}

\begin{tcbpractice}{Pilih 3 aplikasi yang sering digunakan (seperti sosial, belanja, alat), masuk ke halaman manajemen izin aplikasi di [Pengaturan], audit satu per satu, matikan izin sensitif yang tidak perlu atau ubah menjadi "Izinkan hanya saat digunakan".}
Bab 5: Membuat Ponsel Lebih Memahami Anda — Pengaturan Dasar

\end{tcbpractice}
